\begin{tikzpicture}[font=\small,>=stealth,line join=round]
\draw[draw=teal!45!black,fill=green!3,rounded corners=3pt] (0,0) rectangle (13.6,5.9);
\node at (6.8,5.3500000000000005) {中心的纵坐标是 $f(a)$，不一定是 $0$};
\begin{scope}[shift={(1.6,1.1)},x=1.25cm,y=0.35cm]
\draw[gray,->] (-0.3,0)--(3.3,0) node[right] {$x$};
\draw[gray,->] (0,-0.3)--(0,11.5);
\draw[teal!75!black,very thick,domain=-0.2:2.2,samples=100] plot(\x,{(\x-1)^3+7});
\fill[teal!75!black] (0,6) circle(2.5pt);
\fill[teal!75!black] (2,8) circle(2.5pt);
\fill[blue!70!black] (1,7) circle(2.5pt);
\draw[dotted,gray] (0,6)--(2,8);
\draw[red!70!black,thick] (.92,-.25)--(1.08,.25) (.92,.25)--(1.08,-.25);
\node[below] at (1,-.4) {$(1,0)$};
\node[below right] at (1.1,6.5) {$(1,7)$};\node[left] at (0,6) {$6$};\node[above right] at (2,8) {$8$};
\end{scope}
\node[anchor=west,align=left,text width=6cm] at (6.6,2.5) {$f(x)=(x-1)^3+7$。\\[10pt]两点的平均高度为 $(6+8)/2=7$。\\[10pt]蓝点 $(1,7)$ 才是中心；\\红叉 $(1,0)$ 不是。\\[10pt]由 $f^{\prime}$ 的轴只能读出中心横坐标，\\纵坐标还须看 $f(1)$。};
\end{tikzpicture}
